\section{Resumen}
Este artículo presenta un análisis documental acerca de la integración de paradigmas, metodologías y arquitecturas que fortalecen el desarrollo de software contemporáneo. Entre los conceptos revisados se encuentran la Programación Orientada a Objetos (POO) \cite{poo}, el patrón Modelo–Vista–Controlador (MVC) \cite{mvc}, los frameworks Node.js y Spring Boot \cite{nodejs_springboot}, las metodologías ágiles como Scrum \cite{scrum}, y el papel de las bases de datos en los sistemas modernos \cite{databases}. Se empleó un enfoque cualitativo de investigación documental, revisando artículos académicos y aplicando sus hallazgos a un proyecto formativo. Los resultados evidencian que la adecuada articulación entre arquitecturas, metodologías y paradigmas aporta organización, mantenibilidad y escalabilidad a los proyectos. Asimismo, metodologías ágiles como Scrum [3] fortalecen la comunicación, el trabajo colaborativo y la entrega incremental de valor.

Muchas de las investigaciones consultadas fueron aplicadas directamente en el proyecto, como el uso de Scrum \cite{scrum}, que permitió agilizar el desarrollo mediante la organización del equipo de trabajo en Sprints, facilitando el orden de tareas y la obtención de resultados funcionales en cada iteración. Este marco de trabajo resultó especialmente útil en un entorno donde se esperaban avances significativos en tiempos reducidos.

En cuanto a las bases de datos, se identificó la importancia de modelos robustos y escalables que puedan adaptarse a la evolución de las aplicaciones. El uso de sistemas NoSQL se plantea como una alternativa en proyectos que, como el nuestro, requieren integrar el Frontend con una API capaz de gestionar consultas y mantener la coherencia de los datos \cite{databases}.

El patrón Modelo–Vista–Controlador (MVC) \cite{mvc} también se consolidó como una herramienta clave, ya que permite separar las responsabilidades del software: el Modelo (datos y lógica de negocio), la Vista (interfaz de usuario) y el Controlador (flujo de eventos y coordinación). Esta división favorece la mantenibilidad y la escalabilidad del sistema.

Por su parte, la Programación Orientada a Objetos (POO) \cite{poo} fue aplicada como base en el desarrollo con Java, garantizando un nivel adecuado de modularidad, reutilización y adaptabilidad del software. Sus principios —como la herencia, el encapsulamiento y el polimorfismo— resultaron esenciales para mantener la coherencia y facilitar modificaciones posteriores.

Palabras clave: Programación Orientada a Objetos, MVC, Scrum, Node.js, Spring Boot y Bases de datos.
